\documentclass[10pt,a4paper]{amsart}
\usepackage[margin=19mm]{geometry}
\usepackage{amsmath,amssymb,mathtools}
\usepackage[scr=rsfs,cal=euler]{mathalfa}
\usepackage{xfrac}
\usepackage[unicode,hidelinks,bookmarks=false]{hyperref}
\newcommand{\ie}{i.e.\ }
\newcommand{\cf}{cf.\ }
\newcommand{\resp}{respectively\ }
\newcommand{\Subsection}{\subsection}
\providecommand{\nobreakdash}{\nobreak\hbox{-}}
\setlength{\emergencystretch}{2em}
\multlinegap0pt
\usepackage{amssymb}
\usepackage{mathtools}
\usepackage[shortlabels]{enumitem}
\newtheorem{theorem}{Theorem}
\newcommand{\Aut}{\operatorname{Aut}}
\newcommand{\MonG}{\mathbb C B^{\mathbb C^\times}(G)}
\newcommand{\defl}{\operatorname{Def}}
\newcommand{\kerlin}{\mathbb S_p}
\title{The $p$-Biset Functor of Monomial Burnside Rings I: Composition Factors}
\author{\.Ibrahim Kaan Aslan, Olcay Co\c{s}kun}
\date{Source: arXiv:2505.15150v2 (CC BY 4.0)}
\begin{document}
\maketitle

\noindent\emph{Source and licence.} The $p$-Biset Functor of Monomial Burnside Rings I: Composition Factors. Version arXiv:2505.15150v2 (CC BY 4.0), distributed by arXiv under the Creative Commons Attribution 4.0 International licence, http://creativecommons.org/licenses/by/4.0/, as recorded in the arXiv licence metadata for this exact version and shown on the abstract page https://arxiv.org/abs/2505.15150v2. Full licence notice: \texttt{licenses/arxiv-2505.15150-CC-BY-4.0.txt}.\par\smallskip

\noindent\textit{Editorial scope.} Counted unit: the article's theorem thm:cyclic (source0.tex lines 749--765). The article's complete original proof of it is delivered with it (source0.tex lines 775--854, one \texttt{proof} environment of 80 lines, which the article places away from the statement). No mathematics is written, repaired or re-derived: the statement and the proof are the article's own lines, in the article's own notation, and no retained line is substituted or re-flowed.  No further source-local context is needed: every cross-reference of every retained line resolves inside the two delivered spans.  Nothing was omitted from any retained span, so the package carries no omission record.  No retained line contains a citation, so the wrapper carries no bibliography and no bibliography entry.  No retained line contains a figure environment, an image-inclusion command or drawing code, so no image is delivered and the package declares no asset dependency.  Editorial formatting only, in addition: an AMS \texttt{amsart} preamble that restates the article's own declarations and macros used by the retained lines, namely the environments \texttt{theorem} on separate counters, together with the standard packages those lines need.  The retained environments print in this document as Theorem 1 in order of appearance, whereas the article prints the same items with the numbering of its own full document.  The complete native source of the article as distributed in the arXiv e-print for this exact version is bundled unchanged as \texttt{sources/178483/source0.tex}.
\begin{theorem}\label{thm:cyclic}
Let $G = C_{p^k}$ with $k\ge 2$. The set $\mathcal{D}_G$ is a $\mathbb C$-basis for 
$\widetilde{e_G^G}\kerlin(G)$, on which $\Aut(G)$ acts monomially via
\begin{equation}\label{eqn:phistar}
\beta\cdot e_{G, l, \phi} = \phi(\beta|_{H_l})\, e_{G, l, \phi}
\end{equation}
for $\beta\in\Aut(G)$ and $e_{G,l,\phi}\in\mathcal{D}_G$. Moreover there is an isomorphism 
of $\mathbb C[\Aut(G)]$-modules
\begin{equation}\label{eqn:cyclicfinal}
\mathcal K\kerlin(G) \cong \bigoplus_{\substack{2\le l< k\\ \phi_p\,\mathrm{faithful}}}
\mathbb C[e_{G, l, \phi}]
\oplus \bigoplus_{\phi\not = 1}\mathbb C[e_{G, 1, \phi}]
\oplus\, \mathbb C[e_{G, 1, 1} - (1-p)e_{G, 0, 1}],
\end{equation}
and $\mathcal{K}\kerlin(G)$ is isomorphic to the direct sum of all non-primitive irreducible
$\mathbb C[\Aut(G)]$-modules, each with multiplicity $1$.
\end{theorem}

\begin{proof}{[of Theorem ~\ref{thm:cyclic}.]}
The proof proceeds in four steps. We first show that $\mathcal{D}_G$ is a $\mathbb C$-basis for $\widetilde{e_G^G}\kerlin(G)$ and determine the $\Aut(G)$-action on it. We then compute the deflation formula for the basis elements $e_{G,l,\phi}$, from which the isomorphism (\ref{eqn:cyclicfinal}) is read off directly. Finally we verify that the summands afford exactly the non-primitive irreducible characters of $\Aut(G)$, each with multiplicity $1$.

\textbf{Step 1.} The dimension of $\widetilde{e_G^G} \kerlin(G)$ is equal to the number of non-generators of $G$ and it is also easy to see that the cardinality of $\mathcal{D}_G$ is also equal to this number. Thus it is sufficient to check that the set is linearly independent.

If $l\not = n$, then for any choices of $\phi\in\widehat{\Aut(H_l)}$ and $\psi\in\widehat{\Aut(H_n)}$, the primitive idempotents that appear in $e_{G, l, \phi}$ and $e_{G, n, \psi}$ are distinct. Hence there can be no linear relation between these type of elements. Hence the only possible relations are among those elements for fixed $l$. On the other hand, for a fixed $l$, the coefficient matrix for $e_{G, l, \phi}$ as $\phi$ runs over all elements of $\widehat{\Aut(H_l)}$ is equal to the character table of $\widehat{\Aut(H_l)}$, and hence is invertible. In particular the set $\{ e_{G, l, \phi} \mid \phi\in\widehat{\Aut(H_l)}\}$ is linearly independent, which proves the first claim.

For the second claim, the group $\Aut(G)$ acts on the primitive idempotents of $\MonG$ via
\[
\beta\cdot e_{H, h}^G = e_{\beta(H), \beta(h)}^G
\]
for $\beta\in\Aut(G)$ and $e_{H, h}^G\in \MonG$. This action gives rise to a monomial action on $\mathcal{D}_G$. We evaluate this action explicitly: Let $e_{G, l, \phi}\in \mathcal{D}_G$. Then
\[
\beta\cdot e_{G, l, \phi} = \sum_{\alpha\in\Aut(H_l)} \phi(\alpha^{-1})\cdot e_{G, \beta(\alpha(h_l))}^G
\]
Since $G$ is a cyclic group, the restriction of $\beta$ to $H_l$ is an automorphism of $H_l$ and hence $\beta\circ \alpha$ runs over all automorphisms of $\Aut(H_l)$ as $\alpha$ does. Therefore putting $\gamma = \beta\circ \alpha$ we get
\begin{eqnarray*}
\beta\cdot e_{G, l, \phi} &=& \sum_{\gamma\in\Aut(H_l)} \phi(\gamma^{-1}\circ \beta)\cdot e_{G, \gamma(h_l)}^G\\
&=& \phi(\beta|_{H_l}) \cdot e_{G, l, \phi}.
\end{eqnarray*}
Here the last equality holds since $\phi$ is a group homomorphism. Therefore the line $\mathbb C[e_{G, l, \phi}]$ in $\kerlin(G)$ is a $\mathbb C[\Aut(G)]$-submodule with character given by
\begin{equation}
\Aut(G)\to \mathbb C^\times,\ \beta\mapsto \phi(\beta|_{H_l}).
\end{equation}

\smallskip
\noindent\textbf{Step 2.} By the transitivity of deflation maps, it is sufficient to determine the kernel of $\defl^G_{G/Z}$ where $Z:= H_1$ is the unique subgroup of $G$ of order $p$. Also note that any linear relation between deflated basis elements $\defl^G_{G/Z}e_{G, l, \phi}$ may be reduced to a relation for a fixed $l$. Indeed, for distinct numbers $l\not = n, l\ge 2$, we have, for $\alpha\in \Aut(H_l)$ and $\alpha'\in\Aut(H_n)$, that $\alpha(h_l)Z\not = \alpha'(h_n)Z$. Hence the sets of idempotents that appear in deflation of basis elements for $l$ and for $n$ are disjoint and hence no relation between them is possible. Thus we may work with a fixed index $l$.

We claim that
\begin{equation}\label{eqn:deftoZ}
  \defl^G_{G/Z} e_{G, l, \phi} =\begin{cases}
              0 & \text{if } l\ge 2 \text{ and }\ker\phi\cap \Aut_Z(H_l) = \{1\},\\
            e_{G/Z, l-1, \tilde\phi} & \text{if } l\ge 2 \text{ and }\ker\phi\cap \Aut_Z(H_l) \not = \{1\},\\
                          0 & \text{if } l = 1 \text{ and }\phi \not = 1,\\
            \frac{p-1}{p}e_{G/Z, 0, 1} & \text{if } l = 1 \text{ and } \phi = 1,\\
            \frac{1}{p}e_{G/Z, 0, 1} & \text{if } l = 0.
          \end{cases}
\end{equation}
First we assume $l\ge 2$. Note that since $G$ is cyclic, there is no supplement to $Z$ in $G$ and hence the deflation number $m_{G, g}^Z$ is equal to $1/p$ for any $g\in G$. Hence we have
\[
p\cdot \defl^G_{G/Z} e_{G, l, \phi} = \sum_{\alpha\in\Aut(H_l)}\phi(\alpha^{-1}) e_{G/Z, \alpha(h_l)Z}^{G/Z}.
\]
Using the above notation, $\alpha(h_l)Z$ is equal to $\Psi(\alpha)(h_lZ)$. Moreover for any $\kappa_z\in\Aut_Z(H_l)$, we have $\Psi(\alpha) = \Psi(\alpha\kappa_z)$. Hence the right hand side of the above equality can be rearranged as
\begin{eqnarray*}
&=& \sum_{\alpha\in [\Aut(H_l)/\Aut_Z(H_l)]} \Big(\sum_{\kappa\in\Aut_Z(H_l)}\phi(\alpha^{-1}\kappa^{-1})\Big) e_{G/Z, \Psi(\alpha)(h_lZ)}^{G/Z}
\end{eqnarray*}
where the first sum is over a complete set of left coset representatives of $\Aut_Z(H_l)$ in $\Aut(H_l)$. Using the multiplicativity of $\phi$, the right hand side of the above equation becomes
\begin{eqnarray*}
&=& \Big(\sum_{\kappa\in\Aut_Z(H_l)}\phi(\kappa^{-1})\Big)\sum_{\alpha\in [\Aut(H_l)/\Aut_Z(H_l)]} \phi(\alpha^{-1}) e_{G/Z, \Psi(\alpha)(h_lZ)}^{G/Z}.
\end{eqnarray*}
Notice that the coefficient is zero unless $\phi|_{\Aut_Z(H_l)}$ is the trivial character. Moreover when $\phi|_{\Aut_Z(H_l)}$ is the trivial character, it deflates to the quotient $\Aut(H_l)/\Aut_Z(H_l)$. We denote the image of this deflation under the isomorphism $\Psi$ by $\tilde\phi$. Then the above sum becomes
\begin{eqnarray*}
&=& p\sum_{\alpha\in [\Aut(H_l)/\Aut_Z(H_l)]} \tilde\phi(\alpha^{-1}_Z) e_{G/Z, \alpha_Z(h_lZ)}^{G/Z}\\
&=& p\sum_{\alpha_Z\in \Aut(H_l/Z)} \tilde\phi(\alpha^{-1}_Z) e_{G/Z, \alpha_Z(h_lZ)}^{G/Z}\\
&=& p\cdot e_{G/Z, l-1, \tilde\phi}.
\end{eqnarray*}
This proves the first two lines of Equation (\ref{eqn:deftoZ}). The remaining cases are not difficult to deal with. We leave the details to the reader. When $l=1$, we have $H_1 = Z$ and it is easy to see that the above result holds at this extreme and hence we obtain the third and the fourth lines of Equation (\ref{eqn:deftoZ}). The last line also follows since $e_{G/Z, 0, 1} = e_{G/Z, 1}^{G/Z}$.

\smallskip
\noindent\textbf{Step 3.}
From Equation (\ref{eqn:deftoZ}), the basis elements $e_{G, 1, \phi}$ for $\phi\not=1$ together with the combination $e_{G, 1, 1} - (1-p)\cdot e_{G, 0, 1}$ form a basis for the kernel of $\defl^G_{G/Z}$, and the isomorphism (\ref{eqn:cyclicfinal}) follows.

\smallskip
\noindent\textbf{Step 4.}
We need to show that each non-primitive irreducible character of $\Aut(G)$ appears on the right-hand side of Equation (\ref{eqn:cyclicfinal}) with multiplicity 1. First note that, by Equation \ref{eqn:phistar}, the character of the last term is the trivial character. For each $1\le l\le k$ and $\phi\in\widehat{\Aut(H_l)}$, let $\phi^*$ be the character of $\widehat{\Aut(G)}$ given by Equation (\ref{eqn:phistar}). It is clear that if we fix $l$, then $\phi^* = \psi^*$ if and only if $\phi = \psi$.

Fix a pair $(l, \phi)$. Also let $\alpha_p$ be a generator of the Sylow $p$-subgroup of $\Aut(G)$. Then since the $p$-part of $\phi$ is faithful, the value $\phi^*_p(\alpha_p) = \phi_p(\alpha_p\mid_{H_l})$ must be a primitive $p^{l-1}$-th root of unity. In particular, if $(m, \psi)$ is another pair with $l\not= m$, then we also have $\phi^*\not=\psi^*$. Thus there is no repetition in the right hand side of Equation (\ref{eqn:cyclicfinal}).

Moreover since the order of the restriction map $\Aut(G)\to \Aut(H_l)$, $\beta\mapsto \beta\mid_{H_l}$, is $p^{k-l}$, the induced character $\phi^*$ is non-primitive (since $p^{k-l}$ is a quasi-period for $\phi^*$).

We count the number of terms. If $2\le l< k$, then the number of irreducible characters of $\Aut(H_l)$ with faithful $p$-part is $p^{l-2}(p-1)^2$. Summing over $l$, we get
\begin{align*}
\dim_\mathbb C\Big(\bigoplus_{\substack{2\le l< k\\ \phi\in \widehat{\Aut(H_l)}\\ \phi_p:\text{faithful}}}\mathbb C[e_{G, l, \phi}]\Big)&=\sum_{\substack{2\le l< k\\ \phi\in \widehat{\Aut(H_l)}\\ \phi_p:\text{faithful}}}1 \\ &= \sum_{l=2}^{k-1} p^{l-2}(p-1)^2 = (p-1)(p^{k-2}-1).
\end{align*}
With the contribution of the last two terms, the dimension of the right hand side is
\[
(p-1)(p^{k-2}-1) + (p-2) + 1 = (p-1)p^{k-2},
\]
which is equal to the number of non-primitive irreducible characters of $\Aut(G)$. Hence each non-primitive irreducible $\mathbb C[\Aut(G)]$-module appears as a summand in Equation (\ref{eqn:cyclicfinal}) with multiplicity 1, as required.
\end{proof}

\end{document}
